%% toc: P15 173 173.5,173.6,173.7,173.8
\section{Proofs for the range estimate and its omitted variable}
\label{sec:annex-p15-w02b-p15-073}

Notation is that of
\cref{def:p15-w02b-p15-073-range-estimate,def:p15-w02b-p15-073-lane-model}.

\begin{annexitem}{annex:p15-w02b-p15-073-monotone-box}{thm:p15-w02b-p15-073-monotone-box}{Ranges of monotone models and outward rounding}
\begin{proof}
Let \(x\in B\). Change the coordinates of \(x\) one at a time to those of
\(\underline x\). In a nondecreasing coordinate the value moves from
\(x_i\) down to \(l_i\); in a nonincreasing coordinate it moves from \(x_i\)
up to \(u_i\). Each intermediate point lies in \(B\), because each coordinate
stays in its interval, and in either case the step does not increase \(f\).
After \(k\) steps, \(f(\underline x)\leq f(x)\). Changing the coordinates
toward \(\overline x\) instead, each step does not decrease \(f\), so
\(f(x)\leq f(\overline x)\). Both corners lie in \(B\), so
\(f(\underline x)\) is the minimum and \(f(\overline x)\) the maximum of \(f\)
on \(B\).

The box \(B\) is convex, hence connected. A continuous image of a connected
set is connected, and the connected subsets of \(\mathbb R\) are intervals.
The interval \(f(B)\) contains its minimum and maximum, so it equals
\([f(\underline x),f(\overline x)]\).

If \(m'\leq f(\underline x)\) and \(M'\geq f(\overline x)\), then every value
\(f(x)\) satisfies \(m'\leq f(\underline x)\leq f(x)\leq f(\overline x)\leq M'\).

Finally, \(f(x)>c\) for every \(x\in B\) if and only if the minimum of \(f\)
on \(B\) exceeds \(c\). The minimum is attained at \(\underline x\), so this
holds exactly when \(f(\underline x)>c\).
\end{proof}
\end{annexitem}

\begin{annexitem}{annex:p15-w02b-p15-073-lane-flow}{thm:p15-w02b-p15-073-lane-flow}{Lane flow}
\begin{proof}
(1) Number the cars \(k=0,1,2,\ldots\) from the front, and let the front of
car \(0\) be at position \(a\) at time \(0\). In steady traffic the front of
car \(k\) is at \(a-ks+vt\) at time \(t\). It reaches a fixed point \(p\geq a\)
at time
\[
 t_k=\frac{p-a+ks}{v}.
\]
Successive crossing times differ by exactly \(s/v\). In a time interval of
length \(T\) beginning at or after \(t_0\), the number \(N_T\) of crossings therefore satisfies
\(Tv/s-1\leq N_T\leq Tv/s+1\), and \(N_T/T\to v/s\) as \(T\to\infty\).
The car flow is \(v/s\), and the lane flow is \(n\) times it, \(q=nv/s\).
Since \(s\geq L>0\) and \(n,v>0\), \(q\leq nv/L\).

(2) In the time-gap model \(s=L+v\tau\). Substituting in (1),
\[
 q=\frac{nv}{L+v\tau}=\frac{n}{\tau+L/v}.
\]
Then
\[
 \frac{q_0}{q}=\frac n\tau\cdot\frac{\tau+L/v}{n}=1+\frac{L}{v\tau}=1+z,
\]
and \(q\tau/n=\tau/(\tau+L/v)=1/(1+z)\). Since \(L,v,\tau>0\), \(z>0\), so
\(q_0/q>1\) and \(q<q_0\).

(3) The function \(q_0=n/\tau\) is continuous on the box, increasing in \(n\)
and decreasing in \(\tau\). By \cref{thm:p15-w02b-p15-073-monotone-box} its
values form \([n_-/\tau_+,n_+/\tau_-]\).
The function \(q=n/(\tau+L/v)\) is continuous on the box. Its denominator
is positive, increasing in \(\tau\) and in \(L\), and decreasing in \(v\);
so \(q\) is increasing in \(n\) and \(v\) and decreasing in \(\tau\) and
\(L\). The lemma gives its minimum at \((n_-,\tau_+,L_+,v_-)\) and its maximum
at \((n_+,\tau_-,L_-,v_+)\), which are the stated endpoints.
The ratio \(1+L/(v\tau)\) is continuous, increasing in \(L\) and decreasing in
\(v\) and \(\tau\); the lemma gives the stated range with minimum at
\((L_-,v_+,\tau_+)\) and maximum at \((L_+,v_-,\tau_-)\).

(4) By (1), any spacing rule with \(s\geq L\) gives \(q\leq nv/L\).
With \(n\leq n_+\), \(L\geq L_-\) and \(v>0\),
\[
 q\leq\frac{n_+v}{L_-}
 <\frac{n_+}{L_-}\cdot\frac{L_-n_-}{n_+\tau_+}
 =\frac{n_-}{\tau_+}.
\]
By (3), \(n_-/\tau_+\) is the smallest value of \(q_0\) on the box, so \(q\)
lies below every value of \(q_0\) there.

(5) By (2), \(q/q_0=1/(1+L/(v\tau))\). With \(n,\tau,L\) fixed and
\(v\to\infty\), \(L/(v\tau)\to0\), so \(q/q_0\to1\). This is \(q\sim q_0\) in
the sense of \cref{def:p15-approximation-relations}, with parameter
\(\varepsilon=1/v\to0^+\) and \(q_0\neq0\). Since \(q_0\) is independent of
\(v\), also \(q\to q_0\). Strictness of \(q<q_0\) at each finite \(v\) is (2).
\end{proof}
\end{annexitem}

\begin{annexitem}{annex:p15-w02b-p15-073-invalid-candidate}{thm:p15-w02b-p15-073-invalid-candidate}{Invalid candidates}
\begin{proof}
Under a unit change \(\sigma\) each quantity becomes
\(x_i'=x_ie^{d_i^{\mathsf T}\sigma}\) and the target becomes
\(y'=ye^{d_y^{\mathsf T}\sigma}\).

(1) Since \(C\) does not depend on the units,
\[
 g(x')=C\prod_{i=1}^m\bigl(x_ie^{d_i^{\mathsf T}\sigma}\bigr)^{a_i}
 =Cx^a\exp\Bigl(\Bigl(\sum_{i=1}^m a_id_i\Bigr)^{\mathsf T}\sigma\Bigr)
 =g(x)\,e^{(Da)^{\mathsf T}\sigma}.
\]
If \(y=g(x)\), then
\(y'/g(x')=e^{d_y^{\mathsf T}\sigma}/e^{(Da)^{\mathsf T}\sigma}
=e^{(d_y-Da)^{\mathsf T}\sigma}\). For \(\sigma=t(d_y-Da)\) the exponent is
\(t\lVert d_y-Da\rVert^2\). Because \(d_y-Da\neq0\), the norm is positive, and
for \(t\neq0\) the exponent is nonzero, so the ratio differs from \(1\).

(2) Put \(\delta=d_i-d_j\neq0\). For any unit change,
\[
 \frac{x_i'}{x_j'}=\frac{x_i}{x_j}\,e^{\delta^{\mathsf T}\sigma}.
\]
Given \(c>0\), choose
\(\sigma=\bigl(\log c-\log(x_i/x_j)\bigr)\delta/\lVert\delta\rVert^2\). Then
\(\delta^{\mathsf T}\sigma=\log c-\log(x_i/x_j)\), and the ratio equals \(c\).

(3) Here \(n\) has vector \((0,0)\) and is unchanged, while
\(\tau'=\tau e^{\sigma_T}\), \(L'=Le^{\sigma_L}\) and
\(d_y^{\mathsf T}\sigma=-\sigma_T\). The candidate transforms as
\[
 \tilde g(x')=\frac{n}{\tau e^{\sigma_T}+Le^{\sigma_L}} .
\]
The required factor gives \(e^{-\sigma_T}n/(\tau+L)=n/(\tau e^{\sigma_T}+Le^{\sigma_T})\).
Both expressions have the same positive numerator, so they are equal exactly
when \(Le^{\sigma_L}=Le^{\sigma_T}\). Since \(L>0\) and the exponential is
injective, this holds exactly when \(\sigma_L=\sigma_T\).

Now fix \(\sigma_T=0\). Then \(q'=qe^{-\sigma_T}=q\), so the target does not
change. The map \(\sigma_L\mapsto Le^{\sigma_L}\) is a continuous bijection of
\(\mathbb R\) onto \((0,\infty)\), and \(u\mapsto n/(\tau+u)\) is a continuous,
strictly decreasing map of \((0,\infty)\) onto \((0,n/\tau)\): its values lie
strictly between \(0\) and \(n/\tau\), it tends to \(n/\tau\) as \(u\to0^+\)
and to \(0\) as \(u\to\infty\), and by the intermediate value theorem it
takes every value in between. The composite is therefore onto
\((0,n/\tau)\).
\end{proof}
\end{annexitem}
